\paragraph{The fourth power of the theta key.}
Write \(g=f(p+q)\). On the abelian locus let \(L\) be the rigidified
symmetric principal theta line, \(P=\pi_*L\), and
\(\omega=\det(\operatorname{Lie} A)^\vee\), where the dual applies to
the determinant. The line \(P\) commutes with base change. Indeed its
fibres have one section and vanishing higher cohomology. A local
generator \(s\) cuts out a relative effective theta divisor
\(\Theta\). The canonical evaluation identifies
\(\mathcal O(\Theta)=L\otimes\pi^*P^{-1}\), hence
\(e^*\mathcal O(\Theta)=P^{-1}\). In particular this identification
also works when the origin lies on \(\Theta\).

\begin{proposition}\label{ctwist:key}
There is a functorial unit section on the abelian locus, up to a
sign on each parity, of \(P^{\otimes8}\otimes\omega^{\otimes4}\).
Its scalar \(k(s,v)\) in a theta frame \(s\) and a tangent basis
\(v\) transforms by
\[
 k(a s,Mv)=a^{-8}(\det M)^4 k(s,v).
\]
In the metric on the cover its absolute norm is \((2\pi)^{2g}\).
In a primitive theta frame and a full integral tangent basis it is
a unit at good reduction, including residue characteristic two.
\end{proposition}
\begin{proof}
Apply Moret--Bailly, Theorems 0.2 and 0.5 and Corollary 1.5
\cite{moret-bailly-key}, to the effective principally polarized
symmetric divisor. His line is
\(\omega\otimes(e^*\mathcal O(\Theta))^{-2}\).
The scalar is taken relative to \(s^8\otimes\eta_v^4\), where
\(\eta_v(v_1,\ldots,v_g)=1\). This gives both transformation
exponents and the assertion over an abelian scheme over a DVR.
The multiplier-line description of Candelori, Theorem 5.1 and
Remarks 5.2--5.3 \cite{candelori-key}, is over
\(\mathbb Z[1/2,i]\).

Here is the metric calculation. On the even characteristic chart
take periods \((1,\Omega)\), \(Y=\operatorname{Im}\Omega>0\),
the usual Riemann theta, and
\(v_i=(2\pi i)^{-1}\partial/\partial z_i\). The analytic
trivialization of Theorem 0.5 gives \(|k|=1\) in these frames.
Write \(N_s\) for the Haar mean of the squared rigidified theta
norm. With \(z=x+\Omega y\), integration in \(x\) annihilates the
cross terms, and summation of the diagonal terms gives
\(\int_{\mathbb R^g}\exp(-2\pi y^tYy)\,dy=(\det(2Y))^{-1/2}\).
Put \(N_v=\operatorname{covol}(\Lambda)/
(\pi^g|\det(v)|^2)\) in the same complex coordinates. Then
\[
 |k|N_s^4 N_v^2=(2\pi)^{2g}.
\]
The two exponents just computed give invariance under changes of
theta and tangent frames. A symplectic basis puts every even
characteristic in the chosen chart. For an odd characteristic take
the product of the pair with itself. This has even parity, the
two metrics multiply, and Corollary 1.5 identifies the product
trivializations up to sign. Positive square roots give the claimed
identity in odd parity as well. On an orthonormal cover frame the
Hermitian matrix of the polarization is \(2\xi I\) and the theta
has \(N_s=1\). Unimodularity of the alternating form gives
\((2\xi)^g\operatorname{covol}(\Lambda)=1\), so
\(N_v=(2\pi\xi)^{-g}\). The cover's Lie factor for the four
determinants is \((2\pi\xi)^{-2g}\), giving the stated norm.
\end{proof}

There is a further integral change of frame even on the abelian
locus. Let \(R\) be a test DVR whose fraction field splits \(E\),
\(T\) its free Lie lattice, and \(T_\tau=T\cap(T_K)_\tau\) of
rank \(m_\tau\). If \(\mathfrak D_E\) is the different, then
\[
 0\le {\rm length}_R(T/\textstyle\bigoplus_\tau T_\tau)
 \le\sum_\tau m_\tau {\rm ord}_R(\tau\mathfrak D_E).
\]
Indeed for the projector \(e_\tau\), the trace-dual basis formula
\(e_\tau=\sum_i b_i\otimes\tau(b_i^*)\) gives
\(\tau(\mathfrak D_E)e_\tau T\subset T_\tau\):
the \(b_i\) are an integral basis and \(b_i^*\in\mathfrak D_E^{-1}\).
Projection embeds the displayed quotient in
\(\bigoplus_\tau e_\tau T/T_\tau\). Each summand is generated by
at most \(m_\tau\) elements and is killed by
\(\tau(\mathfrak D_E)R\). In the combined eigentangent basis
the valuation of the key increases by four times this length.
All normalized valuations follow by dividing by the ramification
index of the test DVR. Consequently \(k^m\) is integral at these
tests for every \(m\ge0\), with logarithmic complex cost
\(2gm\log(2\pi)\).
\src{src/Development/AbelianVarietyTheory/Proofs/DivisorProofs/ThetaKeySectionNormComplexPointsEqTwoPiPow.lean}{124}{The normalized cover metric};
\src{src/Development/AbelianVarietyTheory/Proofs/DivisorProofs/ExistsKeyUnitPowTwistingSectionOnComponent.lean}{135}{The power of the key and the separate root section}.

\begin{theorem}[Integral key at the semiabelian tests]\label{ctwist:semi}
For the rigidified graph line of Lemma~\ref{intjet:maximum},
\(k(s,v)\) has valuation zero at every semiabelian DVR test in
a primitive theta frame and a full integral tangent basis. The
different bound above therefore applies to the combined eigentangent
basis at all finite places. On each component \(k^m\) supplies the
weight-\(-4m\) multiplier with cost \(2gm\log(2\pi)\).
\end{theorem}
\begin{proof}
It suffices to compute an order after completion and finite
extension. Use the connected model, and split the semistable
reduction after extension. Write \(A/K\) for the generic fibre.
In an additional finite extension make \(A[8]\) rational. Let
\(\mathcal A'\) be the open subgroup of the Néron model generated
by its identity component and the sections of \(A[2]\).
The identity component is the given connected semistable model
(SGA~7, IX, Proposition 3.2, also used in
\cite[IV, \S3.4]{moret-bailly-pinceaux}); this identification
also applies to its semistable base change. The component group
of \(\mathcal A'\) is killed by two, and
\(\mathcal A'[2]\) is finite flat. For the latter assertion the
kernel of multiplication by two is quasi-finite flat on a
semistable group: on each geometric fibre multiplication factors
on the connected semiabelian variety through isogenies of its
torus and abelian quotient, and flatness follows from the
fibrewise flatness criterion. Its generic fibre is thus
schematically dense. Its support is contained in the union of
the finitely many sections extending the generic two-torsion,
which is closed and finite over the base. The kernel is proper
by the valuative criterion, hence is finite. This argument allows
nonreduced special two-torsion.

Moret--Bailly, II, Corollary 1.2.8 with \(n=2\), gives a symmetric
cubical extension \(L'\) on \(\mathcal A'\). Its restriction to
the identity component is the rigidified cubical line of the
test by uniqueness. It is relatively ample by VI, Proposition
2.1, and \(K((L')^2)=\mathcal A'[2]\), since both are the flat
closures of the generic principal two-torsion. The special
fibre of the connected model is geometrically integral. A
primitive theta section has order zero there: if it vanished
along the fibre, division by a uniformizer would give another
integral section. Its nonzero reduction persists under residue
field extension. Lemma~\ref{intjet:maximum} gives this section
canonical norm at most one on the entire analytification.
Consequently its order along every additional component \(C\)
of \(\mathcal A'\) is nonnegative, by the cubical model formula
\({\rm ord}_C(s)=-\log\|s(x_C)\|\) at its Shilov point
(\cite[Proposition 10.3]{deJongShokrieh}). Smoothness gives
normality, so the vertical inequalities and generic regularity
extend the section over \(\mathcal A'\). Restriction back to the
identity component proves that the two theta lattices agree.
This reasoning also compares a test with additional components
by restriction to the identity component. The Lie lattices agree
under open immersions.

Now apply Moret--Bailly, VIII, 1.0--1.3(i)
\cite{moret-bailly-pinceaux}. The base is a normal trait, the
abelian locus is schematically dense, \(K((L')^2)\) is finite,
and the rank \(d\) of the theta module is one. The asserted
canonical trivialization of \(\Delta^{a(g)}\),
\(\Delta=(\pi_*L')^2\otimes\omega\), therefore applies over
\(\mathbb Z[1/d]=\mathbb Z\). Denote it by \(u\). On the integral
stack of abelian schemes with their effective symmetric principal
theta divisors, \(k^{a(g)}/u^4\) is a global unit. Its square is
one by Moret--Bailly, Proposition 1.4(iii) and Corollary 1.5
\cite{moret-bailly-key}. Here one passes between rigidified lines
and effective divisors by the evaluation construction preceding
Proposition~\ref{ctwist:key}. Compatibility of \(u\) with
restriction to the generic point shows
\(a(g){\rm ord}(k)=4{\rm ord}(u)=0\) in the two agreed lattices.
Neither extracting an integral root nor bounding \(a(g)\) is
required. The eigentangent change and Proposition~\ref{ctwist:key}
now give both estimates. Taking the power on the component uses
the inverse only on the characteristic-zero space; extension by
its component idempotent gives the supported multiplier.
\end{proof}

The distinction between nonnegative order and order zero can be
checked on the Tate elliptic curve with parameter \(Q\) and odd
divisor \([0]\). In its connected model
\(e^*\mathcal O([0])=\omega^{-1}\), hence
\(\Delta^4=\omega^{12}\). The modular discriminant is a
nowhere-zero section of this line on the odd abelian stack.
It agrees up to sign with the key by the unit calculation above.
In the invariant Tate frame its scalar is
\(Q\prod_{n\ge1}(1-Q^n)^{24}\), of order \(v(Q)\).
More generally, on a connected semistable model without the
graph condition one always has \(P'\subseteq P\), with \(P'\)
the enlarged model's theta lattice. If \(P'=aP\), the same
canonical-power argument on \(\mathcal A'\) gives
\({\rm ord}(k(s,v))=8{\rm ord}(a)\ge0\). The odd Tate example
has \({\rm ord}(a)=v(Q)/8\) after the extension. Its toric
characteristic is \(1/2\): in the half-characteristic theta
series the toric exponents are shifted by \(1/2\).
In the graph case Lemma~\ref{intjet:maximum} removes this
index and gives the upper different bound as well.

On a normal toroidal chart with these primitive theta and
invariant differential lattices the same argument applies to
the trait at each generic boundary divisor in characteristic
zero. The key and its inverse are units there in full Lie
frames. Thus its multiplication preserves vanishing at those
boundary divisors. The scalar of this fourth-power section and
the existence of the auxiliary one-key section are distinct.
The latter asks for a fourth root up to a scalar in \(F^\times\)
on each component. Over a characteristic-zero component with
trivial geometric monodromy on the symplectic four-torsion this
follows also from Candelori, Definition 3.2 and Theorem 5.1
\cite{candelori-key}. The multiplier torsor is obtained from
the isometry torsor of the symplectic four-group with its theta
characteristic by pushing forward along a character into \(\mu_4\).
In a geometric marking both this torsor and the torsor of
fourth roots of the key are constant: their two fourth-power
trivializations differ by a constant on the universal complex
stack, by the unit calculation in \cite{moret-bailly-key} in
genus \(g\ge2\). Descent of a geometrically constant \(\mu_4\)-torsor
on a geometrically connected component amounts to a torsor over
\(F\), i.e. to a class in \(F^\times/(F^\times)^4\).
The deep level in the statement has odd \(M\ge3\) and a separate
auxiliary marking \(A\). The odd level alone gives no full
four-torsion marking. For the original even weight spaces the
unit \(k^m\) above is intrinsically defined.
\src{src/Development/AbelianVarietyTheory/Proofs/DivisorProofs/ExistsKeySectionFourthRootPerComponent.lean}{93}{The separate descent up to a component scalar}.

\paragraph{Determinants on a characteristic-zero component.}
Let \(\Lambda\) be the projective \(\mathcal O_E\)-lattice of a
component, with level group contained in its unitary isometries
acting as the identity on \(\Lambda/M\Lambda\), \(M\ge3\).
Every such isometry has determinant one over \(E\). In fact its
determinant \(u\) preserves the invertible ideal
\(\bigwedge_{\mathcal O_E}^{p+q}\Lambda\), is \(1\bmod M\)
(localize to check this on free bases), and satisfies
\(|\tau u|=1\) at every embedding. If \(u\ne1\),
\(M^{[E:\mathbb Q]}\le |N_{E/\mathbb Q}(u-1)|
\le2^{[E:\mathbb Q]}\), an impossibility.
For the split half-marking, fixing \(\Lambda/\mathfrak m\Lambda\)
first gives the congruence modulo \(\mathfrak m\). Taking the
adjoint with respect to the perfect alternating integral pairing
gives the conjugate-ideal congruence for the inverse, and hence
for the isometry itself. Indeed the dual of
\(\mathfrak m\Lambda\) is
\(\overline{\mathfrak m}^{-1}\Lambda\). The two ideals are
coprime with product \(M\mathcal O_E\), giving the asserted
principal congruence.

On such a classical complex component the determinant of the
\(\tau\)-part of the standard Betti local system thus has a global
horizontal generator, including for noncompact signatures. On the
period domain the polarization gives the Hodge exact sequence
\[
 0\longrightarrow\omega_\tau\longrightarrow
 H^1_{\rm dR,\tau}\longrightarrow\omega_{\bar\tau}^{\vee}
 \longrightarrow0.
\]
Dualizing its determinant gives
\(\det\operatorname{Lie}_\tau/\det\operatorname{Lie}_{\bar\tau}\),
the line of the elementary section \(s_\tau\).
The norm of a horizontal volume is constant: in a unitary
orthonormal basis the absolute determinant of the Hermitian form
on the full space is one, independently of its positive and
negative Hodge subspaces. Changing to the cover normalization
multiplies this by a constant fixed on the component.
For a smooth geometrically connected algebraic component over
\(F\) carrying the indicated abelian scheme and splitting of \(E\),
the Gauss--Manin determinant is regular singular. Full faithfulness
for regular singular connections gives an algebraic horizontal
generator over \(\mathbb C\) (Deligne, II, Theorem 5.9
\cite{deligne-reg-sing}). It descends up to a scalar to \(F\):
horizontal sections are the kernel of the connection on global
sections, and both global-section spaces and the kernel commute
with extension of scalars on this quasi-compact separated
component. The descended one-dimensional space generates the
line everywhere by faithful flatness. This proves rational
triviality with its archimedean constant for these components.

There are two further consequences for a fixed classical component.
Take a smooth characteristic-zero toroidal compactification with
unipotent local monodromy and with the Hodge bundles given by the
canonical extension of the Gauss--Manin bundle. The horizontal
determinant and its inverse extend as units in the ratio of
the two extended Hodge determinants. Indeed, near a normal
crossings boundary with parameters \(z_i\), the canonical
extension in a multivalued horizontal basis is obtained by
multiplying by
\(\exp(-\sum_i(\log z_i)\log T_i/(2\pi i))\), where \(T_i\)
are the local monodromies. Each \(\log T_i\) is nilpotent on
the \(\tau\)-part, so this matrix has determinant one. Taking
determinants of the extended Hodge sequence identifies its
volume line with the ratio in question. Thus multiplication
in either direction preserves the subcanonical cusp ideal
\(\mathcal O(-D_\infty)\) on these charts. The extended Hodge
sequence and the boundary ideal in the classical PEL model are
given in \href{https://www-users.math.umn.edu/~kwlan/articles/cpt-ram-ord.pdf}
{Lan, Proposition 1.4.2.1 and Definition 1.4.2.9}.

On any fixed finite-type model of this compactification and line
over an open of \(\operatorname{Spec}\mathcal O_F\), the same
section is an integral generator after deleting finitely many more
places. To see this, choose a finite affine trivializing cover.
The coordinates of \(s\) and \(s^{-1}\) on the generic fibre and
the finitely many gluing identities all descend after inverting a
single nonzero base integer. On this open the order is zero at
every test whose lattice is the pullback of that model line.
Away from \(\Delta_E\), the embedding idempotents of the split
étale algebra \(\mathcal O_E\otimes\mathcal O_{F,w}\) are integral,
so formation of its Lie eigensummands commutes with pullback.
In particular this gives the order for the abelian tests extending
to the fixed PEL chart with its universal abelian scheme. The
exceptional primes produced by this argument depend on the
component, horizontal generator and model. Controlling all the
semiabelian tests and the total arithmetic degree in a varying
family requires the comparisons below.

\paragraph{One vertical order on one component.}
Fix a component \(Z\), a host embedding \(\tau\), and a finite place
\(w\) of its coefficient field \(F\). Orders on test traits are
normalized to restrict to \({\rm ord}_w\) on \(F\). The following
divisor calculation explains what an integral model of the parallel
line has to establish. The special fibre may have several irreducible
components with unequal multiplicities.

\begin{lemma}[Vertical order on a proper model]\label{ctwist:vertical}
Let \(\mathcal X\) be integral, normal, proper and flat over a DVR
\(R\) with uniformizer \(\pi\) and nonempty connected special fibre.
Let \(\mathcal L\) be an invertible sheaf, trivialized by \(s\) on the
entire generic fibre. Suppose some positive power of \(\mathcal L\)
is trivial on the reduced special fibre. Write
\({\rm div}(\pi)=\sum_i m_i C_i\) and
\({\rm div}(s)=\sum_i a_i C_i\). Then all \(a_i/m_i\) are the same
rational number \(\alpha\). For every local test \(R\to V\) with
\(\operatorname{Spec}V\to\mathcal X\), where \(V\) is a valuation
ring extending \(R\), the coordinate of \(s\) in an integral generator
has order \(\alpha\,v(\pi)\). If \(R=\mathcal O_{F,w}\) and one
such test is a DVR of ramification index one, then \(\alpha\in\mathbb Z\).
\end{lemma}
\begin{proof}
Set \(\alpha=\min_i a_i/m_i\). Choose \(M>0\) divisible by its
denominator and by the order of the restricted line in the Picard
group. The rational section \(u=\pi^{-M\alpha}s^M\) has nonnegative
order along every prime divisor and order zero at a minimizing
component. It extends over \(\mathcal X\) by normality, checked in
local frames by the height-one intersection criterion.
In a trivialization of \(\mathcal L^M\) on the reduced special
fibre, its restriction is a global function nonzero at that
component's generic point. Global functions on a proper reduced
connected scheme over a field form a field: their algebra is
finite-dimensional and reduced, and connectedness excludes
nontrivial idempotents. Hence \(u\) is a unit on the reduced fibre,
and therefore generates at every point of the special fibre by
Nakayama. It also generates on the generic fibre. Pulling
\(s^M=\pi^{M\alpha}u\) to \(V\) proves the assertion. At an
unramified DVR, \(v(\pi)=1\) and the coordinate has integral order.
\end{proof}
The two scheme facts used here are
\href{https://stacks.math.columbia.edu/tag/0AVB}{Stacks Project,
Tag 0AVB} and
\href{https://stacks.math.columbia.edu/tag/0366}{Section 33.9}.
If the special fibre has just one irreducible component, the same
proof dispenses with its Picard group: after clearing the single
ratio both \(u\) and \(u^{-1}\) extend by normality. In the case
of a reduced irreducible fibre \(M=1\) suffices. For a model over
a ramified base \(R/\mathcal O_{F,w}\), divide its orders by the
ramification index before comparing tests. The rational choice of
frame alone cannot determine \(\alpha\): for the lattice
\(\ell^a\mathbb Z_\ell\) in \(\mathbb Q_\ell\), the frame \(1\)
has integral coordinate \(\ell^{-a}\).

Here is the specific model comparison needed to apply this reasoning.
On the good deep cover write \(U\) for the framed uniformization.
For each fixed \((Z,\tau,w)\), applying the vertical divisor lemma
requires comparing the test frame with local charts on a normal
projective flat model after a finite level and base change. The
comparison must lift
\emph{every} admissible test, match its frame coordinate up to a
unit, and use the same normalized order on the base uniformizer.
Triviality on the reduced fibre, or an equivalent constancy assertion
in these compatible charts, is essential on reducible fibres.
On the dependency path of the exported theorem the generator on
the component is obtained using resolution and GAGA. The source
then uses the associated period frame to compute the degree.
The finite correction between the two frame conventions at a fixed
\((\tau,w)\) is \(q\) times the logarithmic difference between the
different and \(d_0\) in the prescribed local absolute value. Adding
this fixed correction preserves equality of the readings on that
component. The final change of the global degree follows by summing
the finite corrections. Lemma~\ref{ctwist:vertical} proves the
vertical divisor step from its stated model hypotheses. The splitting
model below computes the ordinary integral lattices directly.
\src{src/Development/AutomorphicSpectrum/Proofs/GeometryProofs/FlatFirstFromResolutionAndGagaStatement.lean}{8}{Generator from resolution and GAGA};
\src{src/Development/AutomorphicSpectrum/Proofs/GeometryProofs/KroneckerCharacterFrameBound.lean}{23}{Passage to the component period frame};
\src{src/Development/AutomorphicSpectrum/Proofs/GeometryProofs/KroneckerCharacterPeriodFormula.lean}{122}{Finite correction in the ordinary metric};
\src{src/Development/AutomorphicSpectrum/Proofs/DivisibilityProofs/FlatDegreeStOfResKGaga.lean}{23}{Passage from the constant degree to the corrected section}.

\paragraph{Computing the common degree at a CM point.}
Write \(A_\chi=\operatorname{Re}\log\Omega(E,\chi)\). Assume the
local equality just discussed, and the constant archimedean norm
in the prescribed metric. One point with everywhere good abelian
reduction now suffices to read the uncorrected degree on \(Z\).
In the original construction one chooses an orthogonal \(E\)-frame
in the \emph{specified} complex uniformizing component, constructs
its integral CM object over some \(K/F\), and identifies its framed
generic fibre with a point of \(Z\). This choice works for every
\(\tau\) of the chosen CM orientation on that \(Z\). Write \(T_+\)
and \(T_-\) for the integral Lie lattice intersected with the
\(\tau\)- and \(\bar\tau\)-eigenspaces there, with the induced
polarization metrics. Their ranks are \(p,q\). The frame law
identifies the uncorrected line degree of \(s_\tau\) with
\(\widehat{\deg} T_+-\widehat{\deg} T_-\), using normalized
determinant degrees. Indeed integral local bases compute its finite
coordinate by a ratio of exterior determinants. The same ratio in
orthonormal frames computes its norm at infinity. Changing rational
bases adds principal divisors of degree zero. Passing to \(K\)
preserves the answer because
\(\sum_{v\mid w}e(v/w)f(v/w)=[K:F]\), and each complex embedding
has \([K:F]\) extensions. This calculation uses the integral
eigensublattices even at primes dividing the different.
\src{src/Development/AbelianVarietyTheory/Proofs/DivisorProofs/ExistsCmGoodReductionPointOnComponent.lean}{57}{The CM point on the prescribed component};
\src{src/Development/AbelianVarietyTheory/Proofs/DivisorProofs/FlatLineDegreeEqEigenLatticeDegreeSub.lean}{20}{Finite and infinite determinant readings}.

The cancellation of the component's indices can be seen explicitly.
Let \(S_\pm\subseteq T_\pm\) be the full-rank images spanned by the
differentials of the \(p\) and \(q\) CM factors of the two types,
\(I_\pm\) their determinant index ideals in \(\mathcal O_K\), and
\(g_\pm\) their respective polarization Gram matrices over
\(\mathcal O_E\). Put \(G_\pm=|N_{E/\mathbb Q}\det g_\pm|\).
The CM period comparison for the images gives
\[
 \widehat{\deg}S_+=pA_\chi/2-\frac{\log G_+}{2[E:\mathbb Q]},
 \qquad
 \widehat{\deg}S_-=qA_\chi/2-\frac{\log G_-}{2[E:\mathbb Q]} .
\]
For this mixed-isogeny calculation suppose that the finite comparison
supplies an ideal \(\mathfrak n\) of \(E\) satisfying, with ideals
extended via \(\tau\) when necessary,
\[
 \mathfrak n\bar{\mathfrak n}=(\det g_+\det g_-),\qquad
 I_-(\mathfrak n\mathcal O_K)=I_+(\tau\det g_-).
\]
The definite specialization of the index comparison occurs on the
exported dependency path. The displayed mixed ideal identities are
the hypotheses of the following cancellation calculation.
Proposition~\ref{cmeig:degree} below gives an evaluating point on
the indefinite family for which both Lie lattices are integral
direct sums and the mixed ideal hypotheses can be dispensed with.
\src{src/Development/AbelianVarietyTheory/Proofs/DivisorProofs/DegreeImagesEigenLinesCmUnitaryShimuraObject.lean}{122}{Period and Gram determinant of each image};
\src{src/Development/AbelianVarietyTheory/Proofs/DivisorProofs/CmDefiniteLieIndexIdealEqNormIdeal.lean}{20}{Index comparison for the definite specialization};
\src{src/Development/AbelianVarietyTheory/Proofs/DivisorProofs/CmObjectDefiniteLieEigenLatticeDegreeDifference.lean}{21}{Its degree consequence}.

Taking norms gives \(N\mathfrak n=(G_+G_-)^{1/2}\) and
\(NI_+/NI_-=(N\mathfrak n/G_-)^{[K:\tau E]}\). The finite index
adds \(\log NI_\pm/[K:\mathbb Q]\) to the normalized degree in
passing from \(S_\pm\) to \(T_\pm\). Consequently
\[
 \widehat{\deg}T_+-\widehat{\deg}T_-=(p-q)A_\chi/2 .
\]
All dependence on the two Gram matrices has cancelled. This
explains why the degree is independent of \(Z,\tau\) at fixed
host and signature, although the individual local orders may vary.
It also explains why using the CM point \emph{before} proving local
constancy would leave the orders at the other tests undetermined.

\paragraph{The finite degree in the parallel twist.}
Use \(j_\tau=\lambda_{\tau,p}-1-4m\), \(m=fr+1\). Here the
one-based index \(p\) is the multiplicity at \(\tau\) in the
algebraic weight. The positive row block in the period chart has
size \(p_0=q\). In the
source the elementary flat section \(s_\tau\) has entry \(-1\)
on the entire \(\tau\)-block, so the required product is
\(\prod_\tau s_\tau^{-j_\tau}\). For each \(w\) the comparison
needed with its integral model is the equality of the orders at
all tests over \(w\), with common value
\[
 \rho_{\tau w}+D_{\tau w},\qquad
 D_{\tau w}=q\bigl({\rm ord}_w(\tau\mathfrak D_E)
                            -{\rm ord}_w(d_0)\bigr),
\]
where \(d_0\) is the host prime and \(\rho_\tau\) has finite
support. The source defines the corrected degree using
\(\rho_\tau\); denote it by \(-c_1\) if it is independent of
\(\tau,Z\). The norm of the different and the product formula
give
\[
 [F:\mathbb Q]^{-1}\sum_w D_{\tau w}\log Nw
      =q(\log {\rm rd}(E)-\log d_0),\qquad
 {\rm rd}(E)=|\Delta_E|^{1/[E:\mathbb Q]} .
\]
Indeed the norm of the extended different is
\(|\Delta_E|^{[F:E]}\), while that of \(d_0\mathcal O_F\) is
\(d_0^{[F:\mathbb Q]}\). The uncorrected degree of the product
in \eqref{bfam:flat} is therefore exactly
\begin{equation}\label{ctwist:degree}
 \delta\left(\prod_\tau s_\tau^{-j_\tau}\right)
   =\bigl(c_1-q(\log{\rm rd}(E)-\log d_0)\bigr)
       \sum_\tau j_\tau .
\end{equation}
If \(j_\tau=0\) for every \(\tau\), the multiplier is \(1\) and its
integral cost is zero before any comparison. If only
\(\sum_\tau j_\tau=0\), cancellation in \eqref{ctwist:degree}
uses equality of the corrected degrees for the embeddings with
nonzero exponents. For general \(j_\tau\), suppose as in the composition
that \(|c_1|\le9rL\), \(\log{\rm rd}(E)\le3L\), and
\(|\log{\rm rd}(E)-\log d_0|\le\log{\rm rd}(E)\).
The last inequality follows when \(E\) contains the host quadratic
field of absolute discriminant \(d_0\): the discriminant tower
gives \(2\log{\rm rd}(E)\ge\log d_0\).
At each \(w\) above a prime outside the bad set \(B\), a test of
ramification index one makes the common order integral. Take the
product first, then apply \eqref{bfam:flat} over \(F\) once.
Writing \(W=\sum_{\tau,i}|\lambda_{\tau i}|\), the logarithmic
cost for both twists together is at most
\[
 2gm\log(2\pi)+12rL\,\bigl(W+f(1+4m)\bigr)
 +\sum_{\ell\in B}\log\ell
 +\frac{\log|\Delta_F|}{2[F:\mathbb Q]} .
\]
Here \(q\le r\) and \(\sum_\tau|j_\tau|\le W+f(1+4m)\).
The estimate holds for every component and embedding at the same
threshold as the field, level and degree bounds. On the selected
classical family Theorem~\ref{cmeig:splitmodel} gives the
uncorrected degree \((p-q)A_\chi/2\) of \(s_\tau\). Subtracting
the degree of \(D_\tau\) gives the corrected degree
\[
        -c_1=\tfrac12(p-q)A_\chi-q(\log{\rm rd}(E)-\log d_0).
\]
Thus the product coefficient in \eqref{ctwist:degree} is
\(-\tfrac12(p-q)A_\chi\), including zero in balanced signature.
On the chosen family
the bound on \(c_1\) follows numerically from
\(|A_\chi|\le3L\), \(0\le\log d_0\le L\),
\(0\le\log{\rm rd}(E)\le2L\), and \(p+q=r\).
In fact these bounds give \( |c_1|\le(3/2+2)rL\).
\src{src/Development/AbelianVarietyTheory/Proofs/DivisorProofs/FlatDegreeLogFormAbsLe.lean}{55}{Period and discriminant bounds on the selected stage};
\src{src/Development/AbelianVarietyTheory/Definitions/Foundations/ExistsFlatSecRankZero.lean}{751}{The corrected degree in Lie multiplicities};
\src{src/Development/AbelianVarietyTheory/Proofs/DivisorProofs/ExistsTwSecOfConst.lean}{87}{Conversion to the degree of the actual product};
\src{src/Development/AutomorphicSpectrum/Proofs/GeometryProofs/KroneckerCharacterFrameBound.lean}{20}{The period in the degree of the uncorrected frame};
\src{src/Development/AutomorphicArithmetic/ThetaCorrespondence/ThetaPairLifts/HodgeBundleWeights/ShortSpanCompositionClosed.lean}{54}{The algebraic weight index in the parallel exponent}.

There is a weaker sufficient test when only a lower bound for the
product degree is sought. On a common normal proper flat model at
\(w\) carrying the actual lines and tests of Lemma~\ref{ctwist:vertical},
let \(a_{\tau i}\) be the order of \(s_\tau\) along \(C_i\), and use
the base \(\mathcal O_{F,w}\). Set
\[
 b_w(j)=\min_i\frac{-\sum_\tau j_\tau a_{\tau i}}{m_i}.
\]
Then \(b_w(j)\) is a lower bound at \emph{every} test for this
product without any special-fibre triviality assumption. Indeed,
raise the product to a power \(M\) clearing the denominator of
\(b_w(j)\), and divide by \(\pi^{Mb_w(j)}\). Its divisor is
effective, so normality extends it over the model. Pullback to
each test yields the inequality. A generic trait of a minimizing
vertical divisor realizes equality among all trait maps. On a
reduced special fibre \(m_i=1\) and \(b_w(j)\) is integral.
Thus, if these lower bounds have finite support, are integral
outside \(B\), and give a sufficiently large normalized degree
in \eqref{bfam:flat} with the product's actual metric, the same
integral spanning estimate follows. This criterion applies on each
chosen model separately; a bound uniform in the window still has
to be supplied for the indicated weighted minima.

\paragraph{A product CM point on each indefinite marked row.}
For the fields of Proposition~\ref{host:art} there is another
choice of evaluating point. Write \(k=\mathbb Q(\sqrt{-d_0})\),
\(F_0=E^+\), \(O=O_E\), and \(\Phi\) for the embeddings inducing
the chosen embedding of \(k\). For \(\tau\in\Phi\) write \(p,q\)
for the Lie multiplicities at \(\tau,\bar\tau\), respectively.
Write \(p_0\) for the number of positive Hermitian row directions.
The selected period map conjugates the first \(p_0=q\) coordinates
and is \(\tau\)-linear on the last \(r-p_0=p\) coordinates.
Keep the period convention \(\Pi=\mathfrak a L\) of
Lemma~\ref{bfam:minimum}, so \(L\) is Hermitian self-dual.
\src{src/Formalism/LocalToGlobal/ProjectiveLattices/PetersonNormSubspace.lean}{67}{The conjugated positive block in the period map};
\src{src/Formalism/LocalToGlobal/ProjectiveLattices/PetersonNormSubspace.lean}{75}{The positive row index on a good chart};
\src{src/ReciprocityTheory/ReductiveGroups/ArithmeticGeometry/BirationalGeometry/ToroidalCompactification/FlagBasis.lean}{76}{The Lie multiplicities of the two embeddings}.

\begin{lemma}\label{cmeig:row}
If \(p,q\ge1\), every marked row in this family has, up to rational
isometry, underlying lattice
\[
        J e_1\ \perp\ O e_2\ \perp\cdots\perp\ O e_r,\qquad
        J\bar J=O,\qquad
        h(e_i,e_i)=
        \begin{cases}1&i\le p_0,\\-1&i>p_0.\end{cases}
\]
The transported mark need not split into marks of the summands.
\end{lemma}
\begin{proof}
The two imaginary quadratic subfields of discriminants
\(-d_0,-q_0\) in Proposition~\ref{host:art} have coprime
discriminants (here \(q_0\) is the auxiliary prime). At each
rational prime at least one of them is unramified. Its inertia
restriction is trivial, whereas complex conjugation acts on it
nontrivially. Thus \(E/F_0\) is split or unramified at every finite
place, since relative inertia is the intersection with
\(\operatorname{Gal}(E/F_0)\).

At a split place a basis of the first factor and its pairing-dual
in the second identify a self-dual lattice with the standard one.
At an unramified quadratic place its reduction is a nondegenerate
Hermitian space over the quadratic residue extension. There is a
vector of nonzero norm: otherwise polarization of \(h(x+ay,x+ay)\)
and the perfect residue trace pairing would give \(h=0\).
Lifting gives a unit norm, which can be scaled to either sign,
since the norm maps unramified units onto units (use the residue
norm and lift successively with the residue trace). Its integral
orthogonal complement is again self-dual. This gives the local
orthogonal basis by induction, including in residue characteristic
two. At infinity the signatures agree with the displayed form.
For clarity, the rational isometry follows already from
Hasse--Minkowski: the underlying quadratic form \(h(x,x)\)
represents the prescribed first sign at every place, hence globally.
For the latter implication apply isotropy to \(h(x,x)-a t^2\);
if \(t=0\) gives a nonzero isotropic vector, the original quadratic
space contains a hyperbolic plane and represents \(a\).
Split off the Hermitian line of norm \(a\); by Witt cancellation its
complement is locally isometric to the remaining diagonal form.
Induction now identifies the two rational Hermitian spaces.

Put \(L_0=\bigoplus O e_i\), \(G=U(h)\). Choose
\(g\in G(\mathbb A_{F_0,f})\) carrying \(\widehat L_0\) onto
\(\widehat L\) using the local isometries, choosing the identity
where the two lattices coincide. Write \(t=\det g\) and
\(d(t)=\operatorname{diag}(t,1,\dots,1)\). The determinant-one
groups have strong approximation away from the infinite places:
\(\mathrm{SU}(h)\) is simply connected, absolutely almost simple
(\(r\ge2\)), and noncompact at a real place by \(p,q\ge1\).
This is \href{https://www.mathnet.ru/links/02c4d3528af06a248862fb1aa3ba53f9/im2220_eng.pdf}
{Platonov, Theorem A}. Thus the open set
\(g K_0 d(t)^{-1}\cap\mathrm{SU}(h)(\mathbb A_{F_0,f})\)
contains a rational element \(\gamma\): it contains
\(g d(t)^{-1}\), and \(K_0\) is the open stabilizer of
\(\widehat L_0\). Then \(\gamma^{-1}L\) has completion
\(d(t)\widehat L_0\). Its first ideal is
\(J=E\cap t\widehat O\); \(t\bar t=1\) gives \(J\bar J=O\).
Transport the entire given finite marking by \(\gamma^{-1}\).
The resulting polarized marked row is isomorphic to the one
started with, as in Lemma~\ref{munf:marks}. No approximation to
the identity of the marking, or change of its level, is needed.
\end{proof}

\begin{proposition}[The degree on the product point]\label{cmeig:degree}
On each classical marked period component of
Lemma~\ref{cmeig:row} there is a point over a finite extension of
the coefficient field with everywhere good abelian reduction. For
every \(\tau\in\Phi\) its two integral Lie determinants in the
polarization metric of the cover satisfy
\[
 \widehat{\deg}T_+-\widehat{\deg}T_-=
   \frac{p-q}{2}\left(\gamma_{\rm Eul}+
     \log\frac{2\pi}{{\rm rd}(E)}
       -\frac{L'(1,\chi)}{L(1,\chi)}\right).
\]
All finite places, including those ramified in \(E\), occur with
their ordinary integral eigensublattices.
\end{proposition}
\begin{proof}
In the representative of Lemma~\ref{cmeig:row} take the flag of
the positive and negative coordinate axes. The period lattice
is the orthogonal direct sum of the fractional ideals
\(I_1=\mathfrak a J,\ I_i=\mathfrak a\ (i>1)\).
The last \(p\) factors have CM type \(\Phi\), and the first \(q\)
have its conjugate, by the period map above.
Let \(C\) be the elliptic curve with complex lattice \(O_k\),
with the indicated action (conjugate the action for the opposite
type). It is defined with its endomorphisms over a number field
and acquires good reduction everywhere after finite extension:
its CM \(j\)-value is an algebraic integer, hence gives potential
good reduction at the finitely many places of bad reduction.
The Serre construction \(I_i\otimes_{O_k} C\) exists integrally
since \(I_i\) is projective over \(O_k\): an idempotent on a
finite direct power of \(C\) realizes it as a direct factor.
Its periods at the chosen embedding are exactly the images of
\(I_i\) at the extensions of that \(k\)-embedding. The alternating
trace polarization on each factor is induced from a positive
Hermitian form \(\operatorname{Tr}_{E/k}(\beta u\bar v)\).
Indeed divide its purely imaginary trace coefficient (with the
coordinate sign) by the coefficient of the standard principal
elliptic polarization, a signed generator of
\(\mathfrak D_{k/\mathbb Q}^{-1}\); the ratio \(\beta\in E^+\)
is totally positive by the Riemann positivity condition.
Perfection is exactly
\[
       (\beta)\,I_i\bar I_i\,\mathfrak D_{E/k}=O .
\]
This also follows by taking the trace-dual ideal. Hence the
polarized product has exactly the required period lattice and
polarization. For the polarized tensor construction see
\href{https://arxiv.org/abs/1507.07607}{Amir-Khosravi, \S2}.
The transported finite torsion marking is defined after another
finite extension, irrespective of how it mixes factors. It is
needed on the characteristic-zero point; the underlying product
extends over the full ring of integers. In a classical PEL moduli
chart this constructs the point of the prescribed component,
including its canonical rigidified graph line. An additional
finite étale marking cover is handled by lifting the point to
the given component of the cover, whose map to the underlying
connected component is finite étale and surjective.

Here is the integral and metric calculation for a single factor
with ideal \(I\) and tensor coefficient \(\beta\). Enlarge the
number field \(K\) of definition to split \(E\), and put
\(\mathfrak D=\mathfrak D_{E/k}\). In
\(I\otimes_{O_k}O_K\), the intersection with the coordinate line
at an active embedding \(\nu\) is \(\nu(I\mathfrak D)O_K\).
In fact test membership in the lattice by its dual. The trace
dual of \(I\) over \(O_k\) is \(I^{-1}\mathfrak D^{-1}\);
this remains the dual after flat base change. On the split
generic algebra the trace is the sum of coordinates. The vector
supported at \(\nu\) with entry \(z\) thus pairs integrally
exactly when \(z\nu(I^{-1}\mathfrak D^{-1})\subset O_K\), as
asserted. Tensoring with the invertible \(O_K\)-module
\(\operatorname{Lie} C\) gives the actual Lie eigenline of the
abelian scheme, since taking Lie commutes with the idempotent
defining the tensor construction.

On that coordinate its squared metric at \(\sigma:K\to\mathbb C\)
is the squared elliptic Lie metric multiplied by
\(\sigma\nu(\beta)\): diagonalizing the trace Hermitian form in
the embedding coordinates gives coefficients \(\sigma\nu(\beta)\);
the polarization convention multiplies both Hermitian forms by
the same \(\pi\). For \(e=[E:\mathbb Q]\), normalized ideal and
metric degrees therefore give
\[
\begin{split}
 2\widehat{\deg}T_\nu(I)
 &=2\widehat{\deg}\operatorname{Lie}C
   -\frac{2\log N(I\mathfrak D)+\log|N_{E/\mathbb Q}\beta|}{e}\\
 &=2\widehat{\deg}\operatorname{Lie}C-\frac{\log N\mathfrak D}{e}.
\end{split}
\]
The second equality uses perfection, so even the fractional ideal
on the first coordinate disappears. The norm is the absolute
ideal norm and \(N\mathfrak D=|\Delta_E|/d_0^{e/2}\).

For \(C(\mathbb C)=\mathbb C/\Pi_C\) with Euclidean period area
\(a_C\), principal polarization gives the Hermitian coefficient
\(1/a_C\). The tangent metric is \(\pi/a_C\), so its dual on
\(dz\) is \(a_C/\pi=(i/2\pi)\int dz\wedge d\bar z\).
Thus \(-\widehat{\deg}\operatorname{Lie}C\) is exactly the
stable elliptic Faltings height with the normalization of
\href{https://people.math.wisc.edu/~tonghaiyang/KRY2.pdf}
{Kudla--Rapoport--Yang, (10.14)--(10.15), Proposition 10.10}.
Their quadratic CM formula gives
\[
 2\widehat{\deg}\operatorname{Lie}C
 =\gamma_{\rm Eul}+\log(2\pi/\sqrt{d_0})
                  - L'(1,\chi)/L(1,\chi).
\]
This applies to either elliptic orientation. Substitution gives
the same \(2\widehat{\deg}T_\nu(I)=A_\chi\) on every active
coordinate; orthogonality and the integral direct product give
the two determinants by addition over \(p,q\) coordinates.
\end{proof}

In particular no comparison between a mixed isogeny image and
its saturation is needed to evaluate this point. If the vertical
orders on a selected component agree at all tests and its
archimedean norm is constant, the degree is that of this point
by the normalized base-change calculation preceding
\eqref{ctwist:degree}.

\paragraph{Vertical orders at hyperspecial places.}
There is a direct good-reduction comparison for the ordinary Lie
frames of the classical PEL problem. It also distinguishes a prime
in the level from a prime ramified in the endomorphism algebra.

\begin{proposition}\label{cmeig:hyperspecial}
Suppose the characteristic-zero cover on the chosen component
is a classical finite marking level for its principally polarized
\(O_E\)-abelian scheme, and its finite tests use the integral Lie
eigensublattices of connected semiabelian extensions of that same
scheme. Let \(w\mid\ell\) in the splitting coefficient field \(F\),
with \(\ell\nmid 2\Delta_E\). Each rational horizontal determinant
generator has the same normalized order at all semiabelian DVR
tests on the component above \(w\). If its marked level itself is
hyperspecial at \(\ell\), this common order belongs to \(\mathbb Z\).
\end{proposition}
\begin{proof}
First take a neat prime-to-\(\ell\) marking with full hyperspecial
level at \(\ell\). The algebra is unramified at \(\ell\), its
involution is of unitary type and the principal alternation is
perfect. Thus this is a good reduction PEL problem. With smooth
projective cone data it has a proper smooth toroidal model with
universal semiabelian scheme by
\href{https://www.kwlan.org/articles/cpt-PEL-type-thesis-revision-dense.pdf}
{Lan, Theorems 6.4.1.1 and 7.3.3.4}.
Base change to the completed coefficient DVR. The geometric
components of a proper smooth family form a finite étale scheme;
the selected geometrically connected generic component defined
over the coefficient field extends to an open and closed component
with geometrically connected special fibre. That fibre is smooth,
so is reduced and irreducible. The split embedding idempotents of
\(O_E\otimes_{\mathbb Z}O_{F,w}\) are integral. The indicated Lie
determinant ratio is therefore a line bundle on the model.
Its characteristic-zero horizontal generator extends as a unit at
the toroidal boundary by the unipotent determinant calculation
above Lemma~\ref{ctwist:vertical}. On the whole model its divisor
is an integral multiple of the single reduced special fibre.

Every test extends to the toroidal model by properness. The
pullback of its connected semiabelian group is the given connected
semiabelian extension by uniqueness over a trait (SGA~7 I, IX,
Proposition 3.2). This is an isomorphism respecting the endomorphisms
by generic-fibre uniqueness. Hence its Lie eigensummands, including
those of a toric degeneration, are exactly the test lattices.
Pulling back the divisor in an integral local generator gives the
same normalized integer at every such test.

For general finite marking choose another full neat level of order
prime to \(\ell\) and a geometric connected component of the common
refinement lying over the original component. Over a finite field
extension it and both projection maps are defined. Fix one place
of this field above \(w\). The common component is finite étale
and surjective over each of its two characteristic-zero image
components, as follows by forgetting finite torsion bases. Its two
pulled-back determinant generators differ by a scalar in the
field of definition: both are horizontal for the same connection,
so their quotient has zero differential on the smooth geometrically
connected component. On the auxiliary prime-to-\(\ell\) component
the generator exists by determinant-one monodromy at its full neat
level, as proved above Lemma~\ref{ctwist:vertical}. Any test on the
original component first extends to the fixed base place and
then lifts generically to the common component after a finite
extension of its fraction field, by surjectivity. Connected
semiabelian pullbacks agree by the same uniqueness. The order on
the lifted test is the order on the auxiliary component plus the
valuation of the fixed scalar, both independent of the test after
normalization to \(F_w\). This proves the claim without requiring
an integral level-\(\ell\) marking of a semiabelian group.
\end{proof}

\begin{theorem}[Integral determinant comparison at the ramified primes]
\label{cmeig:splitmodel}
For the fields of Proposition~\ref{host:art}, under the same
classical moduli and test identifications as in
Proposition~\ref{cmeig:hyperspecial}, the common normalized order
exists at every finite place. For each indefinite component the
uncorrected arithmetic degree of the elementary horizontal
determinant is therefore \((p-q)A_\chi/2\).
\end{theorem}
\begin{proof}
Fix \(w,\tau\), and use a neat marking away from its residue
characteristic \(\ell\), with maximal level at \(\ell\) stabilizing
the single self-dual lattice (and its periodic translates).
As in Proposition~\ref{cmeig:hyperspecial} we can pass to this
level and to a finite local splitting field. At every place of
\(E^+\) above \(\ell\) the quadratic algebra is split or
unramified by Lemma~\ref{cmeig:row}. These are exactly the
unitary cases (AL) and (AU) of the Pappas--Rapoport splitting
model. Its order \(O_E\) is maximal, and its alternating lattice
has index one in its dual. These are the hypotheses at the
maximal level of \href{https://stephane-bijakowski.perso.math.cnrs.fr/PRgeometry.pdf}
{Bijakowski--Hernandez, \S2.4 and Theorem 2.30}, which gives a
smooth splitting model in these cases. At two the imaginary
quadratic subfield \(k\) already splits by Lemma~\ref{host:different},
so only the linear case (AL) occurs there.

Use the ordering of the \(E\)-embeddings in
\href{https://www.kwlan.org/articles/cpt-spl.pdf}
{Lan, Choices 2.3.1, Definition 2.3.3}: put \(\tau,\bar\tau\)
first as a consecutive pair in their \(E^+\)-place, and order the
other conjugate pairs consecutively. In the split case the pair
is in distinct integral field factors. In the inert unramified
quadratic case its restrictions to the maximal unramified subfield
are distinct, so the two are again separated by integral
idempotents over the splitting base. On these two different
summands the first quotient of the decreasing Hodge filtration
is labelled respectively by \(\tau,\bar\tau\).
In Lan's interleaved ordering the intervening graded piece on
the other summand is zero. In the increasing notation of
Bijakowski--Hernandez these are the two \emph{last} quotients.
Reversing the ordering within each list, with conjugate lists
paired term by term, is permitted in their \S2.4.

Here are the model comparisons used with the smoothness theorem.
For the single periodic self-dual lattice the open splitting model
parametrizes the polarized Hodge flags of the given ranks. Separating
the unramified idempotents turns Lan's interleaved flags into
the flags for the separate unramified embeddings; conversely they
are interleaved by direct sum. The polarization conditions can
also be compared on the common generic fibre, without identifying
their scheme structures on nonflat thickenings. Indeed both flag
problems are closed in the same relative product of flag varieties
over the naive moduli problem, by the endomorphism and pairing
conditions. They give the same unique flags on the generic fibre.
The Bijakowski--Hernandez model is smooth, hence equals the flat
closure of this generic fibre in that flag space, and is normal.
Lan's normalized open model is the normalization of that closure
on the selected generic components: the forgetful map of the
normalized underlying moduli to the naive moduli is finite and
generically open and closed at full lattice-stabilizer level
(Proposition 2.2.11). Finiteness persists after pullback to the
flag space and after normalization over the excellent base
(Definition 2.4.5 and Proposition 2.3.10). Consequently the normalized
open model on those components is the smooth one, by finite
birationality over a normal scheme.

Lan's Theorem 3.4.1 and Proposition 3.4.14 give a proper flat
toroidal compactification with geometrically regular fibres for
smooth cone data: geometric regularity of each fibre is
equivalent to that of the open splitting model. It is thus smooth.
By Proposition 3.1.2 and Definition 3.1.3 the extended filtration
is a filtration by subbundles on the Hodge bundle of the
semiabelian extension (identified with that of its dual by the
principal polarization, with conjugate labels if using the
dual-action convention). In particular the two locally free
first quotients above exist along the whole boundary.

For any trait write \(H\) for its integral invariant
differentials in the relevant unramified summand. If \(H'\) is
the kernel of its first quotient, then \(H'\) is saturated and
its generic fibre is the sum of all the other eigenspaces in
that summand. Hence
\[
   (H/H')^\vee=\{u\in H^\vee:u(H')=0\}
           =H^\vee\cap (H^\vee_{\rm frac})_\tau .
\]
The identical argument on the conjugate summand gives the
second test lattice. This uses the end quotient of the
filtration: in general an intermediate graded piece has a
different lattice. The determinant ratio of these dual quotients
is therefore an invertible sheaf which gives exactly the
ordinary test metric at \(\ell\). Its generic horizontal section
is a unit across the characteristic-zero boundary, as before.
Over the fixed geometric component the smooth proper model
has a single reduced special fibre. Lemma~\ref{ctwist:vertical}
and the common-refinement argument of
Proposition~\ref{cmeig:hyperspecial} give equality of all orders,
including on semiabelian tests. A different ordering or model
can be used for a different pair \(\tau,\bar\tau\), since orders
are calculated pairwise before multiplying sections.

Finally compare with the CM point of Proposition~\ref{cmeig:degree}
at each place and with its archimedean norm. Its finite extension
preserves normalized arithmetic degrees. This gives the claimed
global degree and hence the uncorrected product degree in
\eqref{ctwist:degree}.
\end{proof}

\paragraph{The selected good cover.}
The level condition on \(Y\) fixes its auxiliary marking to have
order one. Its finite marking is the half-marking at the odd integer
\(M\ge3\).
The characteristic-zero comparison can be made from the moduli
property. The universal object is proper and semiabelian, hence
abelian. Its rigidified line is relatively ample and its
polarization kernel is trivial, hence its polarization \(\lambda\)
is principal. The adjunction property for its biextension gives
the specified Rosati involution on \(O_E\). Conversely from these
PEL data pull the normalized Poincar\'e bundle back along
\((1,\lambda w_0)\). Its polarization is
\(\lambda(w_0+\bar w_0)=\lambda\); the rigidification, symmetry
and cube law are the canonical ones. Its biextension is the
one for \(\lambda\). This reconstruction also recovers the
original line from its graph condition. The rigidified isomorphism
is unique: an automorphism is multiplication by a unit on the
abelian scheme, which comes from the base by properness and
geometric connectedness and equals one at the identity.

Over the splitting field the determinant condition is precisely
the \(p,q\) multiplicity condition. On a reduced
characteristic-zero base the Hodge eigensummands are vector bundles
and their ranks can be tested at geometric points. The mark kills
\(\mathfrak m\) and on the complementary summand is the reduction
of a polarized Tate marking. By CRT these restrictions determine
it on the whole lattice. Its compatibility at all torsion orders
in the definition of an object is a condition on the local
lattice genus, not a choice of additional sections. More explicitly,
in a connected complex PEL family the Betti lattice, its pairing,
action and finite mark are locally constant; a compatible marked
Tate isometry at one point transports to every point. Equivalently
at a geometric point the reduction of a compatible tower is a
point of the prescribed compact-open marking orbit. Conversely
any representative of that orbit gives the required compatible
reductions at all \(n\). To check the roots of unity, fix the
primitive \(M\)-th root in the coefficient field and, at a geometric
point, extend it to a compatible system \((\zeta'_n)\). Relative
to the tower \((\zeta_n)\) of a Tate marking there is a unique
\(u\in\widehat{\mathbb Z}^{\times}\) with \(\zeta'_n=\zeta_n^u\)
for all \(n\ge1\). The reference lattice has an
integral similitude of multiplier \(u\) fixing the half-mark.
At \(\ell\mid M\), use \(1\) on the marked half and \(u_\ell\) on
its pairing-dual half. At \(\ell\nmid M\), use a central scalar of
norm \(u_\ell\) at each place of \(F_0\) over \(\ell\), using the
surjectivity of split and unramified unit norms in Lemma~\ref{cmeig:row}.
Precomposing a Tate marking by this similitude gives the prescribed
tower with the same marked sections. At order \(M\) the Weil pairing
with the marked basis and the fixed root then determines the basis
on the opposite half uniquely. This construction works over the
base at order \(M\) by perfect duality of finite étale groups. The
higher towers occur only in the geometric-point existence test.
This identifies the half-marking with the corresponding fine PEL
level on the selected components. After finite extension the
source moduli property gives a classifying map for the classical
universal object on its reduced smooth base, and the classical
moduli property gives the reverse map on the reduced smooth
source base. The reconstructed universal objects are isomorphic,
so both composites are the identity by uniqueness. The maps
preserve the underlying group and polarization, also after
restricting to a specified marked row.
\src{src/ReciprocityTheory/ReductiveGroups/ArithmeticGeometry/BirationalGeometry/ToroidalCompactification/TwoWeightEqOneAddRoot.lean}{131}{The compatible torsion isometries and primitive roots};
\src{src/ReciprocityTheory/ReductiveGroups/ArithmeticGeometry/BirationalGeometry/ToroidalCompactification/TwoWeightEqOneAddRoot.lean}{144}{The two marks, signature and geometric-point existence condition};
\src{src/ReciprocityTheory/ReductiveGroups/ArithmeticGeometry/BirationalGeometry/ToroidalCompactification/PushInvertible.lean}{29}{The universal object and its reduced-base moduli property};
\src{src/Development/AbelianVarietyTheory/Definitions/Divisors/NormNegCmConstantEqAtAllEmbeddings.lean}{187}{The selected level, period charts and polarization scale}.

On a test DVR the definition supplies a semiabelian group whose
generic fibre is cartesian over this universal group. Its fibres
are geometrically connected. The test vectors form a basis for
all integral eigenvectors at each labelled embedding, so they give
exactly the intersections used in Theorem~\ref{cmeig:splitmodel}.
Indeed the images of algebraic integers in the coefficient field
belong to the valuation ring; the equality expressing the
eigenvector property therefore holds integrally if and only if
it holds over the fraction field, by torsion-freeness of the
smooth Lie lattice. The theta generator in the test contributes
no factor to the parallel character. At infinity the comparison
uses the polarization metric \(\pi H\) on Lie vectors; this is
the metric used in Proposition~\ref{cmeig:degree}. Thus the
ordinary parallel determinant comparison follows on the selected
principal family with its given period readings.
\src{src/Formalism/LocalToGlobal/ProjectiveLattices/VAtVNeBot.lean}{53}{Cartesian semiabelian tests and the universal integral eigenbasis condition};
\src{src/ReciprocityTheory/ReductiveGroups/ArithmeticGeometry/SemiabelianSchemes/AbelianSchemeQuotients/IsAbelianScheme.lean}{41}{Connected geometric fibres of the semiabelian group};
\src{src/Development/AbelianVarietyTheory/Definitions/CMLattices/IsDiffOf.lean}{101}{Polarization norm in the CM comparison}.

\begin{corollary}[Short sections in the final weight on a good cover]
\label{cmeig:final}
Fix one good deep cover \(Y\) as in Corollary~\ref{crev:good}
on a selected host, with \(p,q\ge1\). Take the admissible final
weight \(\lambda\) and the level, denominator, field and bad-prime
size bounds of Theorem~\ref{bfam:oncharts}.
On each marked component \(Z\), its regular cusp rule space of
weight \(\lambda\) is spanned over the coefficient field by
integral sections \(s_i\) with
\[
     \log\|s_i\|_{\sigma}\le
          C\bigl((frL)^3 W+(frL)^5\bigr),\qquad
     W=\sum_{\tau,i}|\lambda_{\tau i}|,
\]
at every embedding in the theorem. The same constant works over
these choices under the common stage bounds.
\end{corollary}
\begin{proof}
Start with the integral generators of Corollary~\ref{crev:good}.
Multiply by the key of Theorem~\ref{ctwist:semi} and by the
product with exponents \(-j_\tau\) in \eqref{ctwist:degree}.
Theorem~\ref{cmeig:splitmodel} computes the latter degree at the
actual integral eigenlattices. Outside the excluded primes the
level is hyperspecial, so the common order is integral by
Proposition~\ref{cmeig:hyperspecial}. At the remaining primes
round the product order once as in \eqref{bfam:flat}.
Proposition~\ref{host:period} supplies the uniform numerical
degree bound. Both unit sections and their inverses preserve the
cusp condition by Corollary~\ref{crev:twist}. The final-weight
estimate of Theorem~\ref{bfam:oncharts} now gives the claim.
This argument starts with \(Y\) and the stated numerical bounds.
Selecting such covers compatibly at all stages and identifying the
restriction maps on the modified flags enter separately in the
Hilbert--Samuel estimate.
\end{proof}

